\fancyhead[RO]{Introduzione}
\chapter*{Introduzione}

% \section*{Abstract}

% Questa tesi di Laurea triennale in Matematica si pone come obiettivo quello di dimostrare un ben noto risultato della geometria differenziale (ovvero il seguente teorema) tramite gli strumenti forniti dall'analisi non-standard. 

% \paragraph{Teorema} Se $ X $ e $ Y $ sono due campi vettoriali di classe $ C^{2} $ su una varietà differenziabile $ M $, i loro flussi commutano se e solo se $ [X,Y] \equiv 0 $.

% In particolare, nei capitoli 2 e 3 si fornisce una rigorosa trattazione della costruzione dell'estensione non-standard di $ \R $ ed una generalizzazione all'universo $ V(\R) $. Nel capitolo 4, infine, si utilizzano le basi teoriche poste in precedenza per definire l'estensione non standard di una varietà differenziabile $ M $, e per definire la naturale traduzione dei campi vettoriali in ambito infinitesimale. 

% Con queste premesse, il flusso di un campo vettoriale, inizialmente definito come oggetto ``integrale'', diventa un cammino composto da passi infinitesimali. Questo consente una definizione formale di flusso che concordi con la intuizione naturale del concetto.

\section*{Cenni storici}

Quando Lebniz sviluppò per la prima volta gli oggetti dell'analisi matematica (quali funzioni continue, derivate, integrali\dots) il suo lavoro si fondava sul concetto intuitivo di infinitesimo. Per lui gli infinitesimi erano quei numeri diversi da zero ma minori di ogni numero reale positivo.

Leibniz supponeva che per questi elementi valessero le normali regole aritmetiche di calcolo. Vi erano però delle contraddizioni logiche, come messo in luce da George Berkeley. 

Pertanto, nel diciannovesimo secolo vi fu la necessità di riformare le basi dell'analisi matematica; questa necessità culminò nel 1861, con l'introduione da parte di Weierstrass della definizione $ (\varepsilon$-$\delta) $ di limite come la conosciamo oggi. 

Questo portò le dimostrazioni ad essere molto più complesse e poco intuitive. 

Solo con Abraham Robinson fu possibile dare una solida base teorica all'utilizzo degli infinitesimi nell'analisi matematica:  nel 1961 e nel 1966 pubblicò due lavori, rispettivamente \cite{robinson_non-standard_1961} e \cite{robinson_non-standard_1966}, che presentano in modo solido e rigoroso una teoria dei numeri infinitesimali.

Egli scoprì infatti che $ \R $ poteva essere esteso ad un campo non archimedeo\footnote{
    Un campo ordinato si dice archimedeo se per ogni $ x \in \K $ positivo, esiste $ n \in \N $ tale che\[x<\parentesi{n\text{-volte}}{1+\dots+1}.\]
}$\starR$ (detto \emph{campo degli iperreali}) in cui sono quindi presenti degli elementi infiniti e degli elementi infinitesimi, il cui comportamento è quello supposto da Leibniz.

\section*{Campo degli iperreali}

Il campo ordinato dei numeri reali si immerge elementarmente nel campo ordinato degli iperreali. Questo implica anche che ogni enunciato esprimibile al prim'ordine vero per i numeri reali è vero anche per gli iperreali. 

Il Capitolo 2 è dedicato proprio alla costruzione di $ \starR $. Questo si costruisce ponendo una relazione di equivalenza su $ \R^{\N} $ (l'insieme delle successioni a valori reali), sfruttando la teoria degli ultraprodotti. 

In particolare, se $ \mathcal{U} $ è un ultrafiltro non principale su $ \N $, ovvero $ \mathcal{U} \subseteq \parti{\N} $ tale che \begin{itemize}
    \item se $ A \in \mathcal{U} $ e $ A \subseteq B \subseteq \N $, allora $ B \in \mathcal{U} $;
    \item se $ A,B \in \mathcal{U} $, allora $ A \cap B \in \mathcal{F} $; 
    \item $ \emptyset \notin \mathcal{U} $;
    \item per ogni $ A \subseteq \N $, si ha che $ A \in \mathcal{U} $ oppure $ \N \setminus A \in \mathcal{U} $
\end{itemize}si definisce la relazione di equivalenza $ \sim_{\mathcal{U}}  $ tra elementi di $ \R^{\N} $. 

Per ogni \mbox{$ \sequence{a_{n}}, \sequence{b_{n} } \in \R^{\N} $}, \begin{equation*}
    \sequence{a_{n} }\mathrel{\sim_{\mathcal{U}} }\sequence{b_{n} }\quad\text{ se e solo se }\quad\set{ n \in \N: a_{n}=b_{n}  } \in \mathcal{U}.
\end{equation*}Si ha quindi che $ \starR $ è il quoziente di $ \R^{\N} $ con $ \sim_{\mathcal{U}}  $. Si scriverà: \begin{equation*}
    \starR \coloneqq \R^{\N}/\mathcal{U},
\end{equation*}e si indicheranno con $ \filtersequence{a_{n}} $ gli elementi di $ \R^{\N}/\mathcal{U} $. 

Questa costruzione di $ \starR $ prende il nome di \emph{ultrapotenza} di $ \R $. È possibile sfruttare un potente teorema della teoria degli ultrapodotti e delle ultrapotenze: il Teorema di Łoś. Questo consente di affermare che esista una immersione elementare da \begin{equation*}
    \left(\R,+,-,\cdot , 0,1,\le \right)\quad\text{a}\quad\left(\starR,+,-,\cdot , 0,1,\le \right)
\end{equation*}ovvero che esista una funzione \begin{equation*}
    *: \R \longrightarrow \starR
\end{equation*}tale che, se $ \mathcal{L}_{\text{ord}} $ è il linguaggio del prim'ordine dei campi ordinati, per ogni $ \mathcal{L}_{\text{ord}} $-formula $ \phi(x_1,\dots,x_{n} ) $ e per ogni $ a_1,\dots,a_{n} \in \R  $: \begin{equation*}
    \R\vDash\phi[a_1,\dots,a_{n} ]\quad\text{se e solo se}\quad \starR\vDash\phi[\nonstandard{a_1},\dots,\nonstandard{a_{n} }].
\end{equation*}In particolare, questa iniezione è \begin{align*}
* :\R  &\hookrightarrow\starR \\
r &\longmapsto \filtersequence{r}.
\end{align*}Siccome è possibile identificare ogni numero reale con un elemento di $ \starR $, per ogni $ r \in \R $ si scriverà $ r $ in luogo di $ \nonstandard{r} \in \starR $. 

All'interno di $ \starR $ somma, prodotto e relazione d'ordine vengono definite in maniera canonica; ad esempio, per ogni $ \filtersequence{a_{n}}, \filtersequence{b_{n}} \in \starR$: \begin{equation*}
    \filtersequence{a_{n}} + \filtersequence{b_{n}} = \filtersequence{a_{n} + b_{n} }.
\end{equation*}
%  \begin{itemize}
%     \item $ \displaystyle \filtersequence{a_{n}} + \filtersequence{b_{n}} = \filtersequence{a_{n} + b_{n} } $;
%     \item $ \displaystyle \filtersequence{a_{n}} \cdot  \filtersequence{b_{n}} = \filtersequence{a_{n} \cdot  b_{n} }$;
%     \item $ \displaystyle \filtersequence{a_{n}} \le \filtersequence{b_{n}} $ se e solo se $ \set{ n\in \N: a_{n}\le b_{n}  } \in \mathcal{U} $;
% \end{itemize}

Si verifica che l'elemento $ M \in \starR $: \begin{equation*}
    M \coloneqq \filtersequence{n}
\end{equation*}sia \emph{infinito}, ovvero per ogni $ r \in \R $, $ M > r $. 

Ovviamente quindi $ 1/M $ sarà un elemento infinitesimo di $ \starR $, ovvero per ogni $ r \in \R$ positivo \begin{equation*}
    0 < \abs{\frac{1}{M}} <r.
\end{equation*}Il campo $ \starR $ è dunque corredato di infiniti ed infinitesimi: un numero iperreale $ x \in \starR $ si dice \begin{itemize}
    \item \emph{finito} se esiste un numero reale $ r \in \R $ tale per cui $ |x|<r $;
    \item \emph{infinito} se per ogni numero reale $ r \in \R$, $ r<|x| $;
    \item \emph{infinitesimo} se per ogni numero reale positivo $ r \in \R $, $ r>0 $, vale \begin{equation*}
        0 < |x| < r.
    \end{equation*}
\end{itemize}
Poiché per ogni $ x \in \starR $ finito esiste un unico $ r_{x}  \in \R $ tale che $ x-r_{x}  $ è infinitesimo o nullo, è ben definita la funzione $ \st: \starR_{\bounded}  \longrightarrow \R $, che a ciascun $ x $ associa $ r_{x}  $, dove $ \starR_{\bounded} \subseteq \starR $ sono gli elementi finiti di $ \starR $.

È possibile estendere questa costruzione ad un insieme molto più grande di oggetti: l'universo matematico. Vi sono diversi approcci. In questa tesi si seguirà quello di \cite{nowik_differential_2015}, che a sua volta di rifà a \cite{chang_model_2012}.

Sia $ X $ un insieme base\footnote{La definizione formale di questo concetto sarà riportata nel corpo di questo lavoro. È possibile dimostrare che $ \R $ è un insieme base.}. Siano \begin{align*}
    V_0(X) &= X, \\\forall\, n \in \N\setminus\{0\}\quad V_{n+1}(X) &= V_{n}(X)\cup\parti{V_{n}(X) }.
\end{align*}Si definisce \emph{universo matematico} costruito su $ X $ l'insieme \begin{equation*}
    V(X)\coloneqq \bigcup_{n \in \N} V_{n}(X). 
\end{equation*}
Una estensione non standard di $ V(X) $ consiste in un insieme base $ \nonstandard{X} $, con $ X \subsetneqq \nonstandard{X} $, un sotto universo $ \nonstandard{V(X)} \subseteq V(\nonstandard{X}) $, e una funzione \begin{equation*}
    *: V(X) \longrightarrow V(\nonstandard{X})
\end{equation*}che soddisfi le seguenti condizioni. \begin{enumerate}
    \item[(I)] $ \nonstandard{X} \subseteq \nonstandard{V(X)} $, e $ * $ è una mappa iniettiva tale per cui, se $ x \in X $, allora $ \nonstandard{x} \in \nonstandard{X} $. 
    \item[(II)] Per ogni formula limitata $ \phi(x_1,\dots,x_{k} ) $ e ogni $ k $-upla di costanti $ a_1,\dots,a_{k} \in V(X)  $, la formula $ \phi[a_1,\dots,a_{k} ] $ è vera in $ V(X) $ se e solo se $ \phi[\nonstandard{a_1},\dots,\nonstandard{a_{k} }] $ è vera in $ \nonstandard{V(X)} $.
\end{enumerate}

L'intero Capitolo 3 è dedicato alla dimostrazione dell'esistenza di $ \nonstandard{V(X)} $, con una costruzione simile a quella di $ \starR $.  L'oggetto di interesse per questa tesi è $ V(\R) $, la cui utilità è dovuta al fatto che tra i suoi elementi vi siano: \begin{itemize}
    \item sottoinsiemi di $ \R^{n} $ per $ n \in \N $ arbitrario;
    \item funzioni $ f: A \subseteq \R^{n} \longrightarrow B \subseteq \R^{m} $ per $ n,m \in \N $ arbitrari.
\end{itemize}

\section*{Geometria differenziale}

Se $ M $ è una varietà differenziabile, si può identificare l'insieme dei punti di $ M $ con un elemento di $ V_1(\R)\setminus V_0(\R) $. In questo contesto, quindi, ha senso considerare l'estensione non-standard di tutti gli oggetti di interesse della geometria differenziale, quali atlanti differenziabili, topologia di $ M $ e campi vettoriali e scalari. In questo modo sarà ben definito $ \nonstandard{M} $, con $ M \hookrightarrow \nonstandard{M} $.

L'obiettivo di questa trattazione è la definizione di flusso di un campo vettoriale su una varietà come un oggetto combinatorico invece che integrale. In seguito si mostrerà, come applicazione, la dimostrazione di un corollario del Teorema di Frobenius. 

% Sia dunque $ \lambda \in \starR $ un \emph{infinitesimo} fissato.

% È possibile definire, su $ \nonstandard{M} $, una nozione di distanza infinitesima tra punti (e si dirà che i punti sono infinitamente vicini), ereditata proprio dalla struttura non archimedea di $ \starR $. 

% Questo consente di definire i cosiddetti \emph{prevettori}, ovvero coppie di punti di infinitamente vicini. A meno di una relazione di equivalenza (dettata dal fatto che una distanza può assumere una quantità non numerabile di valori infinitesimi), ogni vettore \emph{classico} $ X $ in un punto $ p $ induce un unico prevettore $ (p,x) $, dove $ x \in \nonstandard{M} $ è infinitamente vicino a $ p $ e, in coordinate, $ x = p + \lambda X $.  

% Un prevettore $ (a,x) $ agisce come derivazione sulle funzioni differenziabili $ f: M \longrightarrow \R $: \begin{equation*}
%     (a,x)\,f \coloneqq \frac{1}{\lambda}\left(\nonstandard{f}(x)-\nonstandard{f}(a)\right)
% \end{equation*}Per come è stata definita la distanza infinitesima tra punti, questo è necessariamente un elemento finito di $ \starR $, e pertanto se ne può considerare la  \begin{equation*}
%     \st\left(\frac{1}{\lambda}\left(\nonstandard{f}(x)-\nonstandard{f}(a)\right)\right).
% \end{equation*}

% È possibile dunque definire un \emph{campo di prevettori} come una mappa $ F: \nonstandard{U} \longrightarrow \nonstandard{V} $ per opportuni intorni $ \nonstandard{U} \subseteq \nonstandard{V} \subseteq \nonstandard{M}$ che a ciascun punto $ a \in \nonstandard{U} $ associa $ F(a) $ infinitamente vicino ad $ a $. 

% Si dirà che un campo di prevettori $ F $\ emph{realizza} un campo vettoriale classico $ X $ se in ogni punto di $ U $ le derivazioni di $ X $ e di $ F $ coincidono

% $ \nonstandard{M} $ in\begin{equation*}
% F:\nonstandard{U}\longrightarrow \nonstandard{V}
% \end{equation*}, definite per opportuni intorni $ \nonstandard{U} \subseteq \nonstandard{V} \subseteq M $, che a ciascun punto $ a $ di $ \nonstandard{U} $ associano un punto $ F(a) $ infinitamente vicino ad $ a $. A meno di una relazione di equivalenza, ogni campo vettoriale 

Il primo passo da compiere è quello di definire una nozione di punti infinitamente vicini su $ M $. Sia dunque $ \lambda \in \starR$ un infinitesimo fissato. 

Siano $ a,b \in \nonstandard{M} $. Si scriverà \begin{itemize}
    \item $ a-b \prec \lambda $ se, in coordinate, \begin{equation*}
        \norma{a-b} = r \lambda
    \end{equation*}per qualche $ r \in \R\setminus\{0\} $;
    \item $ a-b\pprec \lambda $ se, in coordinate, \begin{equation*}
        \norma{a- b} = \mu \lambda
    \end{equation*}per qualche $\mu \in \starR$ infinitesimo.
\end{itemize}In entrambi i casi si dirà che $ a,b $ sono \emph{infinitamente vicini}. Queste relazioni si potranno utilizzare anche sostituendo numeri iperreali $ t \in \starR $ ad $ a-b $, utilizzando $ \abs{t} $ in luogo di $ \norma{a-b} $. 

%% Scrivere una nozione per \prec e \pprec

% Per $ r, s \in \starR $, si adotterà la seguente notazione: \begin{itemize}
%     \item $ r \prec s $ se $ r = a \cdot s $ per qualche $ a $ finito;
%     \item $ r \pprec s $ se $ r = \lambda \cdot s $ per qualche $ \lambda $ infinitesimo.
% \end{itemize}Inoltre, se $ v \in \nonstandard{(\R^{n})} $ si scriverà $ v\prec s $ e $ v\pprec s $ se per qualche norma $ \norma{\cdot } $ su $ \R^{n} $ la relazione vale per $ \nonstandard{\norma{v}} $. Per $ v,w \in \nonstandard{(\R^{n})} $ si scriverà $ v\approx w $ se $ v-w \pprec 1$. 



% Si definisce l'\emph{alone} di $ p $ come $ \halo(p) \subseteq \nonstandard{M} $ l'insieme di tutti gli $ x \in \nonstandard{M} $ per i quali vi è un intorno coordinato $ U $ di $ p $ tale che $ x \in \nonstandard{U} $ e $ x \approx p$ in quelle date coordinate. Per ogni $ A \subseteq M $ è definito l'alone di $ A $: \begin{equation*}
%     \hhalo{A} \coloneqq \bigcup_{p \in A}\halo(p).
% \end{equation*}Questa unione è disgiunta, e pertanto è ben definita una funzione $ \st: \hhalo{A} \longrightarrow A $ che ad $ x $ associa l'unico $ p \in A $ tale per cui $ x \in \halo(p) $.

Inoltre, poiché l'obiettivo finale è di studiare oggetti ``standard'', è conveniente definire $ \hhalo{M} \subseteq \nonstandard{M}$ come l'insieme dei punti di $ \nonstandard{M} $ che sono infinitamente vicini a punti standard di $ M $. La definizione si estende naturalmente a $ \hhalo{U} $ per $ U \subseteq M $.

È ora possibile dunque dare una definizione nonstandard per i vettori tangenti ad $ M $.\begin{itemize}
    \item Per $ a \in \hhalo{M} $ si definisce \emph{prevettore} in $ a $ una coppia $ (a,x) $ tale che, in coordinate, \begin{equation*}
        x-a\pprec \lambda.
    \end{equation*}Si indicherà con $ P_{a}M  $ l'insieme dei prevettori in $ a $. 
    \item Due prevettori $ (a,x), (a,y) $ sono equivalenti se per ogni funzione differenziabile $ f: M \longrightarrow \R $ \begin{equation*}
        \nonstandard{f(y)}-\nonstandard{f(x)}\pprec \lambda.
    \end{equation*}Si denota con $ T_{a}M $ il quoziente di $ P_{a}M  $ rispetto a questa relazione di equivalenza. Il quoziente $ T_{a}M  $ è uno spazio vettoriale reale indipendente dalla scelta delle coordinate, i cui elementi sono denotati con $ \vec{ax} $.
\end{itemize}

Se $ a \in \hhalo{M} $ e $ (a,x) $ è un prevettore ed $ \vec{ax} $ è un vettore, si definiscono le derivazioni su $ f: M \longrightarrow \R $: \begin{align*}
    (a,x)\,f& \coloneqq \frac{1}{\lambda}\left(\nonstandard{f}(x)-\nonstandard{f}(a)\right) \in \starR\\
\vec{ax} \,f &\coloneqq \st\left((a,x)\,f\right) \in \R.
\end{align*}

Ogni vettore \emph{classico} $ X $ in $ a $ induce un unico vettore non standard $ \vec{ax} \in T_{a}M $, dove in coordinate $ x = a+ \lambda X $, nel senso che $ X $ e $ \vec{ax} $ agiscono allo stesso modo come derivazioni.

Un campo di prevettori è una mappa interna \begin{equation*}
    F: \nonstandard{M} \longrightarrow \nonstandard{M}
\end{equation*}tale che per ogni $ a \in \hhalo{M} $ valga $ F(a) -a\prec \lambda$.

Per $ U \subseteq V \subseteq M $, un \emph{campo di prevettori locale} è una mappa interna \begin{equation*}
    F: \nonstandard{U} \longrightarrow \nonstandard{V}
\end{equation*}tale che per ogni $ a \in \hhalo{U} $ valga $ F(a)-a \prec \lambda$. Un campo di prevettori realizza un campo vettoriale classico $ X $ se per ogni funzione differenziabile $ h: U \longrightarrow \R $ si ha \begin{equation*}
    \forall\, a \in U:\quad X\,h(a) = \vec{a\,F(a)}\,h
\end{equation*}

Se $ X $ è un campo vettoriale classico, allora il campo di prevettori scritto in coordinate: \begin{align*}
F:\nonstandard{U} &\longrightarrow \nonstandard{V} \\
a &\longmapsto a + \lambda \cdot X(a)
\end{align*}realizza $ X $.

% Vi è quindi la necessità di dare delle condizioni di regolarità ai campi di prevettori. In particolare $ F: \nonstandard{U} \longrightarrow \nonstandard{V} $ è un campo di prevettori di classe $ D^{k} $ se, per ogni $ a \in \hhalo{U} $ e per ogni $ v_1,\dots,v_{k} \in \nonstandard{(\R^{n})} $ con $ v_{i} \prec \lambda  $ si ha che in coordinate: \begin{equation*}
%     \Delta_{v_1,\dots,v_{k}}^{k} F(a) \prec \lambda\norma{v_1}\dots\norma{v_{k} },
% \end{equation*}dove si utilizza la notazione delle differenze finite.

% Se $ X $ è un campo vettoriale classico di classe $ C^{k} $, allora il corrispondende campo di prevettori $ F(a) = a + \lambda \cdot X(a) $ è di classe $ D^{k} $. 

Si può inoltre dimostrare che se $ F: \nonstandard{U} \longrightarrow \nonstandard{V} $ è un campo di prevettori che soddisfa determinate condizioni di regolarità, a meno di ridimensionare il dominio, la mappa $ F $ ammette un'inversa $ F^{-1}: \nonstandard{U} \longrightarrow \nonstandard{V} $.

Con queste nozioni è possibile definire il flusso di un campo di prevettori globale per tempi positivi. Siano $ F: \nonstandard{M} \longrightarrow \nonstandard{M}  $ un campo di prevettori locale, e $ t>0 $ con $ t \in \starR $ e  $ n = n(t) = \floor{t/\lambda} $. Sia $ \Phi $ la mappa \begin{align*}
    \Phi:\operatorname{Map}(M) \times \N &\longrightarrow \operatorname{Map}(M) \\
    (f,n) &\longmapsto \underset{n\text{ volte}}{\underbrace{f\circ \dots\circ f}},
    \end{align*}dove $ \operatorname{Map}(M) $ è l'insieme delle mappe $ M \longrightarrow M$, e se ne consideri l'estensione non standard $ \nonstandard{\Phi}: \nonstandard{\operatorname{Map}(M)}\times \nonstandard{\N} \longrightarrow \nonstandard{\operatorname{Map}(M)} $. Per ogni $ a \in \nonstandard{M} $ il flusso di $ F $ al tempo $ t $ è definito come $ F_{t}(a) = \nonstandard{\Phi}\left(F,n(t)\right)(a) $. 

In questo modo si è formalizzata l'idea intuitiva di flusso come insieme di spostamenti infinitesimi: per costruzione, ciascuna successiva applicazione di $ F $ genera uno spostamento infinitesimo. Ciò nonostante, siccome $ n(t) $ varia su $ \nonstandard{\N} $ e in particolare per $ t $ reali $ n(t) $ è \emph{infinito}, la composizione di spostamenti infinitesimi diventa diventa una ``composizione infinita'', e pertanto lo spostamento totale risulta apprezzabile. 

Pertanto, ha senso ``tornare'' alla varietà classica $ M $, per mezzo del \emph{flusso standard}.

% È possibile estendere a campi di prevettori locali e per tempi negativi la costruzione di cui sopra: questo è trattato approfonditamente nel Capitolo 4. 

Questo, infatti, si può definire a partire dal flusso di un campo di prevettori $ F $, in un intorno $ B \subseteq M$ adeguato, e per tempi adeguati, come: \begin{align*}
h^{F}_t:B &\longrightarrow M \\
x &\longmapsto \st\left(F_{t}(x) \right).
\end{align*}Questo è un oggetto classico, e si può dimostrare che se $ X $ è un campo vettoriale classico e $ F(a) =a + \lambda \cdot X $, allora $ h_{t}^{F}  $  è esattamente il flusso di $ X $ nel senso classico. 
% Si può considerare l'estensione non standard di $ h_{t}^{F} $, ovvero il campo di prevettori: \begin{equation*}
%     \tilde{F} \coloneqq h_{\lambda}^{F}: \nonstandard{B} \longrightarrow \nonstandard{M}.
% \end{equation*}

Dunque, se $ F $ e $ G $ sono campi di prevettori sufficientemente regolari, si definisce la parentesi di Lie come il campo di prevettori: \begin{equation*}
    [F,G] \coloneqq \left(G^{-1}\circ F^{-1}\circ G\circ F\right)^{\floor{1/\lambda}}.
\end{equation*}Inoltre, è significativo notare che per ogni $ a $ nel dominio di $ [F,G] $: \begin{equation*}
    [F,G](a) - a - \frac{1}{\lambda}\left(G^{-1}\circ F^{-1}\circ G\circ F(a)-a\right)\pprec \lambda.
\end{equation*}Infatti, se $ X $ e $ Y $ sono campi vettoriali classici di classe $ C^{2} $ di flussi $ \Theta^{X}_t  $ e $ \Theta^{Y}_t $, allora la parentesi di Lie tra $ X $ e $ Y $ è, per ogni $ p \in M $: \begin{equation*}
    [X,Y]_{\text{cl}}(p) = \lim_{t\to 0} \frac{1}{t^{2}}\left(\left(\Theta^{Y}_t\right)^{-1}\circ\left(\Theta^{X}_t\right)^{-1}\circ \left(\Theta^{Y}_t\right)\circ \left(\Theta^{X}_t\right)(p)-p\right).
\end{equation*}È naturale dedurre che se $ X,Y $ sono campi vettoriali classici e $ F(a)=a + \lambda X(a) $, $ G(b) = b+\lambda Y(b) $ sono campi di prevettori che realizzano, rispettivamente, $ X $ e $ Y $, allora $ [F,G] $ realizza $ [X,Y]_{\text{cl}} $.

Si termina il lavoro dimostrando, tramite un metodo puramente combinatorico, che per due campi vettoriali classici $ X, Y$ si ha che i loro flussi commutano se e solo se $ [X,Y]_{\text{cl}}=0 $.